\documentclass[10pt,a4paper]{amsart}
\usepackage{amsmath,amssymb,amsthm,graphicx}
\usepackage{verbatim}
\newenvironment{reference}{\comment}{\endcomment}
%\newenvironment{reference}{}{}
\newenvironment{slogan}{\comment}{\endcomment}
\newenvironment{history}{\comment}{\endcomment}

% For commutative diagrams we use Xy-pic
\usepackage[all]{xy}

% We use 2cell for 2-commutative diagrams.
\xyoption{2cell}
% Theorem environments.
%
\theoremstyle{plain}
\newtheorem{theorem}[subsection]{Theorem}
\newtheorem{proposition}[subsection]{Proposition}
\newtheorem{lemma}[subsection]{Lemma}

\theoremstyle{definition}
\newtheorem{definition}[subsection]{Definition}
\newtheorem{example}[subsection]{Example}
\newtheorem{exercise}[subsection]{Exercise}
\newtheorem{situation}[subsection]{Situation}

\theoremstyle{remark}
\newtheorem{remark}[subsection]{Remark}
\newtheorem{remarks}[subsection]{Remarks}

\numberwithin{equation}{subsection}

% Macros
%
\def\lim{\mathop{\mathrm{lim}}\nolimits}
\def\colim{\mathop{\mathrm{colim}}\nolimits}
\def\Spec{\mathop{\mathrm{Spec}}}
\def\Hom{\mathop{\mathrm{Hom}}\nolimits}
\def\Ext{\mathop{\mathrm{Ext}}\nolimits}
\def\SheafHom{\mathop{\mathcal{H}\!\mathit{om}}\nolimits}
\def\SheafExt{\mathop{\mathcal{E}\!\mathit{xt}}\nolimits}
\def\Sch{\mathit{Sch}}
\def\Mor{\mathop{\mathrm{Mor}}\nolimits}
\def\Ob{\mathop{\mathrm{Ob}}\nolimits}
\def\Sh{\mathop{\mathit{Sh}}\nolimits}
\def\NL{\mathop{N\!L}\nolimits}
\def\CH{\mathop{\mathrm{CH}}\nolimits}
\def\proetale{{pro\text{-}\acute{e}tale}}
\def\etale{{\acute{e}tale}}
\def\QCoh{\mathit{QCoh}}
\def\Ker{\mathop{\mathrm{Ker}}}
\def\Im{\mathop{\mathrm{Im}}}
\def\Coker{\mathop{\mathrm{Coker}}}
\def\Coim{\mathop{\mathrm{Coim}}}

% Boxtimes
%
\DeclareMathSymbol{\boxtimes}{\mathbin}{AMSa}{"02}

%
% Macros for moduli stacks/spaces
%
\def\QCohstack{\mathcal{QC}\!\mathit{oh}}
\def\Cohstack{\mathcal{C}\!\mathit{oh}}
\def\Spacesstack{\mathcal{S}\!\mathit{paces}}
\def\Quotfunctor{\mathrm{Quot}}
\def\Hilbfunctor{\mathrm{Hilb}}
\def\Curvesstack{\mathcal{C}\!\mathit{urves}}
\def\Polarizedstack{\mathcal{P}\!\mathit{olarized}}
\def\Complexesstack{\mathcal{C}\!\mathit{omplexes}}
% \Pic is the operator that assigns to X its picard group, usage \Pic(X)
% \Picardstack_{X/B} denotes the Picard stack of X over B
% \Picardfunctor_{X/B} denotes the Picard functor of X over B
\def\Pic{\mathop{\mathrm{Pic}}\nolimits}
\def\Picardstack{\mathcal{P}\!\mathit{ic}}
\def\Picardfunctor{\mathrm{Pic}}
\def\Deformationcategory{\mathcal{D}\!\mathit{ef}}
\usepackage{cleveref}
\crefname{originallemma}{original lemma}{original lemmas}
\Crefname{originallemma}{Original lemma}{Original lemmas}
\makeatletter
\newlabel{lemma-twisted-inverse-image}{{0A9E}{source}{}{sourcepointer-1}{}}
\newlabel{lemma-twisted-inverse-image@cref}{{[originallemma][0][]0A9E}{[1][0][]source}{}{}{}}
\makeatother
\crefname{originalreference}{original reference}{original references}
\Crefname{originalreference}{Original reference}{Original references}
\makeatletter
\newlabel{remark-relative-dualizing-complex}{{0B6S}{source}{}{sourcepointer-2}{}}
\newlabel{remark-relative-dualizing-complex@cref}{{[originalreference][0][]0B6S}{[1][0][]source}{}{}{}}
\makeatother
\crefname{originalsection}{original section}{original sections}
\Crefname{originalsection}{Original section}{Original sections}
\makeatletter
\newlabel{section-upper-shriek}{{0A9Y}{source}{}{sourcepointer-3}{}}
\newlabel{section-upper-shriek@cref}{{[originalsection][0][]0A9Y}{[1][0][]source}{}{}{}}
\makeatother
\crefname{originallemma}{original lemma}{original lemmas}
\Crefname{originallemma}{Original lemma}{Original lemmas}
\makeatletter
\newlabel{lemma-shriek-dualizing}{{0AA3}{source}{}{sourcepointer-4}{}}
\newlabel{lemma-shriek-dualizing@cref}{{[originallemma][0][]0AA3}{[1][0][]source}{}{}{}}
\makeatother
\crefname{originalsection}{original section}{original sections}
\Crefname{originalsection}{Original section}{Original sections}
\makeatletter
\newlabel{example-proper-over-local}{{0AWI}{source}{}{sourcepointer-5}{}}
\newlabel{example-proper-over-local@cref}{{[originalsection][0][]0AWI}{[1][0][]source}{}{}{}}
\makeatother
\crefname{originalsection}{original section}{original sections}
\Crefname{originalsection}{Original section}{Original sections}
\makeatletter
\newlabel{example-equidimensional-over-field}{{0AWJ}{source}{}{sourcepointer-6}{}}
\newlabel{example-equidimensional-over-field@cref}{{[originalsection][0][]0AWJ}{[1][0][]source}{}{}{}}
\makeatother
\crefname{originallemma}{original lemma}{original lemmas}
\Crefname{originallemma}{Original lemma}{Original lemmas}
\makeatletter
\newlabel{lemma-vanishing-good-dualizing}{{0AWN}{source}{}{sourcepointer-7}{}}
\newlabel{lemma-vanishing-good-dualizing@cref}{{[originallemma][0][]0AWN}{[1][0][]source}{}{}{}}
\makeatother
\crefname{originallemma}{original lemma}{original lemmas}
\Crefname{originallemma}{Original lemma}{Original lemmas}
\makeatletter
\newlabel{lemma-good-dualizing-normalized}{{0AWL}{source}{}{sourcepointer-8}{}}
\newlabel{lemma-good-dualizing-normalized@cref}{{[originallemma][0][]0AWL}{[1][0][]source}{}{}{}}
\makeatother
\crefname{originallemma}{original lemma}{original lemmas}
\Crefname{originallemma}{Original lemma}{Original lemmas}
\makeatletter
\newlabel{dualizing-lemma-depth-in-terms-dualizing-complex}{{0AWR}{source}{}{sourcepointer-9}{}}
\newlabel{dualizing-lemma-depth-in-terms-dualizing-complex@cref}{{[originallemma][0][]0AWR}{[1][0][]source}{}{}{}}
\makeatother
\crefname{originallemma}{original lemma}{original lemmas}
\Crefname{originallemma}{Original lemma}{Original lemmas}
\makeatletter
\newlabel{lemma-dualizing-module-proper-over-A}{{0AWP}{source}{}{sourcepointer-10}{}}
\newlabel{lemma-dualizing-module-proper-over-A@cref}{{[originallemma][0][]0AWP}{[1][0][]source}{}{}{}}
\makeatother
\crefname{originallemma}{original lemma}{original lemmas}
\Crefname{originallemma}{Original lemma}{Original lemmas}
\makeatletter
\newlabel{lemma-smooth-proper}{{0BRT}{source}{}{sourcepointer-11}{}}
\newlabel{lemma-smooth-proper@cref}{{[originallemma][0][]0BRT}{[1][0][]source}{}{}{}}
\makeatother
\title{Existence and properties of the dualizing complex of a proper scheme over a field}
\author{The Stacks Project Authors}
\date{stacks-project:duality.tex:local-snapshot-20261012; GNU FDL 1.2 or later}
\begin{document}
\maketitle
\noindent\textit{Editorial note.} All delivered mathematical text is the author's own, copied from the pinned Stacks Project chapter file duality.tex; no statement or proof step was rewritten, reordered or removed.\par\medskip
\noindent\textit{Editorial note.} Preamble provenance: the pinned chapter file opens with an input of the shared preamble and carries no preamble text of its own, so the project's single Stacks-wide preamble staged at sources/stacks-duality/preamble.tex is excerpted verbatim under the recipe fixed by LANE\_BRIEF 22.1, as three spans, lines 8-13 (verbatim package and the reference/slogan/history environments), lines 14-18 (Xy-pic setup) and lines 154-217 (the native theorem declarations and mathematical macros). Nothing was added, removed or redefined; the wrapper class and the amsmath/amssymb/amsthm/graphicx base packages are supplied by the assembler.\par\medskip
\noindent\textit{Editorial note.} The counted statement and proof contain no cross reference, no citation and no included file, so the delivered body is self contained.\par\medskip
\noindent\textit{Editorial note.} Canonical location of the counted proposition: Stacks tag 0FVV (sources/stacks-duality/pinned-original-source.tex line 6934).\par\medskip
\noindent\textit{Allowed standard prerequisites (editorial).}\begin{itemize}
\item The right adjoint of pushforward and the relative dualizing complex. Same chapter, Lemma labelled lemma-twisted-inverse-image and Remark labelled remark-relative-dualizing-complex, cited at proof lines 6965-6972. Defines omega and identifies it with f\textasciicircum{}! of the structure sheaf.
\item Identification of the right adjoint with a dualizing complex. Same chapter, Lemma labelled lemma-shriek-dualizing, cited at proof lines 6974-6976. Makes omega a dualizing complex.
\item Vanishing and support statements for a good dualizing complex, and its normalisation at a point. Same chapter, Lemmas labelled lemma-vanishing-good-dualizing and lemma-good-dualizing-normalized, cited at proof lines 6984-6986 and 6990-6992, and Dualizing Complexes, Lemma labelled dualizing-lemma-depth-in-terms-dualizing-complex, cited at proof lines 6992-6994. Give properties (1) to (4).
\item The dualizing module of a proper scheme and the smooth case. Same chapter, Lemma labelled lemma-dualizing-module-proper-over-A, cited at proof lines 7000-7002, and Lemma labelled lemma-smooth-proper, cited at proof lines 7006-7008. Give properties (6) and (7).
\end{itemize}\medskip
\ifdefined\hypertarget\hypertarget{sourcepointer-1}{}\fi
\noindent\textit{Source reference.} Original lemma 0A9E: sources/stacks-duality/pinned-original-source.tex, source lines 653--653. Same-chapter pointer into the author environment lemma at source line 653; invoked at proof line 6968 of the counted proof. It is retained only as a navigation pointer; every load-bearing step is inside the delivered span.\par\medskip
\ifdefined\hypertarget\hypertarget{sourcepointer-2}{}\fi
\noindent\textit{Source reference.} Original reference 0B6S: sources/stacks-duality/pinned-original-source.tex, source lines 2945--2945. Same-chapter pointer into the author environment remark at source line 2945; invoked at proof line 6973 of the counted proof. It is retained only as a navigation pointer; every load-bearing step is inside the delivered span.\par\medskip
\ifdefined\hypertarget\hypertarget{sourcepointer-3}{}\fi
\noindent\textit{Source reference.} Original section 0A9Y: sources/stacks-duality/pinned-original-source.tex, source lines 3940--3940. Same-chapter pointer into the author environment lemma at source line 3940; invoked at proof line 6976 of the counted proof. It is retained only as a navigation pointer; every load-bearing step is inside the delivered span.\par\medskip
\ifdefined\hypertarget\hypertarget{sourcepointer-4}{}\fi
\noindent\textit{Source reference.} Original lemma 0AA3: sources/stacks-duality/pinned-original-source.tex, source lines 4708--4708. Same-chapter pointer into the author environment lemma at source line 4708; invoked at proof line 6977 of the counted proof. It is retained only as a navigation pointer; every load-bearing step is inside the delivered span.\par\medskip
\ifdefined\hypertarget\hypertarget{sourcepointer-5}{}\fi
\noindent\textit{Source reference.} Original section 0AWI: sources/stacks-duality/pinned-original-source.tex, source lines 6070--6070. Same-chapter pointer into the author environment example at source line 6070; invoked at proof line 6980 of the counted proof. It is retained only as a navigation pointer; every load-bearing step is inside the delivered span.\par\medskip
\ifdefined\hypertarget\hypertarget{sourcepointer-6}{}\fi
\noindent\textit{Source reference.} Original section 0AWJ: sources/stacks-duality/pinned-original-source.tex, source lines 6085--6085. Same-chapter pointer into the author environment example at source line 6085; invoked at proof line 6981 of the counted proof. It is retained only as a navigation pointer; every load-bearing step is inside the delivered span.\par\medskip
\ifdefined\hypertarget\hypertarget{sourcepointer-7}{}\fi
\noindent\textit{Source reference.} Original lemma 0AWN: sources/stacks-duality/pinned-original-source.tex, source lines 6125--6125. Same-chapter pointer into the author environment lemma at source line 6125; invoked at proof line 6985 of the counted proof. It is retained only as a navigation pointer; every load-bearing step is inside the delivered span.\par\medskip
\ifdefined\hypertarget\hypertarget{sourcepointer-8}{}\fi
\noindent\textit{Source reference.} Original lemma 0AWL: sources/stacks-duality/pinned-original-source.tex, source lines 5805--5805. Same-chapter pointer into the author environment lemma at source line 5805; invoked at proof line 6990 of the counted proof. It is retained only as a navigation pointer; every load-bearing step is inside the delivered span.\par\medskip
\ifdefined\hypertarget\hypertarget{sourcepointer-9}{}\fi
\noindent\textit{Source reference.} Original lemma 0AWR: sources/stacks-chow/tags.txt, source lines 13729--13729. Cross-chapter standard prerequisite located in the project tag list staged in the local snapshot; invoked at proof line 6992 of the counted proof. It is retained only as a navigation pointer; every load-bearing step is inside the delivered span.\par\medskip
\ifdefined\hypertarget\hypertarget{sourcepointer-10}{}\fi
\noindent\textit{Source reference.} Original lemma 0AWP: sources/stacks-duality/pinned-original-source.tex, source lines 6157--6157. Same-chapter pointer into the author environment lemma at source line 6157; invoked at proof line 7002 of the counted proof. It is retained only as a navigation pointer; every load-bearing step is inside the delivered span.\par\medskip
\ifdefined\hypertarget\hypertarget{sourcepointer-11}{}\fi
\noindent\textit{Source reference.} Original lemma 0BRT: sources/stacks-duality/pinned-original-source.tex, source lines 3879--3879. Same-chapter pointer into the author environment lemma at source line 3879; invoked at proof line 7007 of the counted proof. It is retained only as a navigation pointer; every load-bearing step is inside the delivered span.\par\medskip
\begin{lemma}
\label{lemma-duality-proper-over-field}
Let $X$ be a proper scheme over a field $k$. There exists a dualizing complex
$\omega_X^\bullet$ with the following properties
\begin{enumerate}
\item $H^i(\omega_X^\bullet)$ is nonzero only for $i \in [-\dim(X), 0]$,
\item $\omega_X = H^{-\dim(X)}(\omega_X^\bullet)$ is a coherent
$(S_2)$-module whose support is the irreducible components
of dimension $\dim(X)$,
\item the dimension of the support of $H^i(\omega_X^\bullet)$ is at most $-i$,
\item for $x \in X$ closed the module
$H^i(\omega_{X, x}^\bullet) \oplus \ldots \oplus H^0(\omega_{X, x}^\bullet)$
is nonzero if and only if $\text{depth}(\mathcal{O}_{X, x}) \leq -i$,
\item for $K \in D_\QCoh(\mathcal{O}_X)$ there are functorial
isomorphisms\footnote{This property
characterizes $\omega_X^\bullet$ in $D_\QCoh(\mathcal{O}_X)$
up to unique isomorphism by the Yoneda lemma. Since $\omega_X^\bullet$
is in $D^b_{\textit{Coh}}(\mathcal{O}_X)$ in fact it suffices to consider
$K \in D^b_{\textit{Coh}}(\mathcal{O}_X)$.}
$$
\Ext^i_X(K, \omega_X^\bullet) = \Hom_k(H^{-i}(X, K), k)
$$
compatible with shifts and distinguished triangles,
\item there are functorial isomorphisms
$\Hom(\mathcal{F}, \omega_X) = \Hom_k(H^{\dim(X)}(X, \mathcal{F}), k)$
for $\mathcal{F}$ quasi-coherent on $X$, and
\item if $X \to \Spec(k)$ is smooth of relative dimension $d$,
then $\omega_X^\bullet \cong \wedge^d\Omega_{X/k}[d]$ and
$\omega_X \cong \wedge^d\Omega_{X/k}$.
\end{enumerate}
\end{lemma}

\begin{proof}
Denote $f : X \to \Spec(k)$ the structure morphism. Let $a$ be the
right adjoint of pushforward of this morphism, see
Lemma \ref{lemma-twisted-inverse-image}. Consider the relative dualizing
complex
$$
\omega_X^\bullet = a(\mathcal{O}_{\Spec(k)})
$$
Compare with Remark \ref{remark-relative-dualizing-complex}.
Since $f$ is proper we have
$f^!(\mathcal{O}_{\Spec(k)}) = a(\mathcal{O}_{\Spec(k)})$ by
definition, see Section \ref{section-upper-shriek}.
Applying Lemma \ref{lemma-shriek-dualizing} we find that
$\omega_X^\bullet$ is a dualizing complex. Moreover, we see that
$\omega_X^\bullet$ and $\omega_X$ are as in
Example \ref{example-proper-over-local} and as in
Example \ref{example-equidimensional-over-field}.

\medskip\noindent
Parts (1), (2), and (3) follow from
Lemma \ref{lemma-vanishing-good-dualizing}.

\medskip\noindent
For a closed point $x \in X$ we see that $\omega_{X, x}^\bullet$ is a
normalized dualizing complex over $\mathcal{O}_{X, x}$, see
Lemma \ref{lemma-good-dualizing-normalized}.
Part (4) then follows from Dualizing Complexes, Lemma
\ref{dualizing-lemma-depth-in-terms-dualizing-complex}.

\medskip\noindent
Part (5) holds by construction as $a$ is the right adjoint to
$Rf_* : D_\QCoh(\mathcal{O}_X) \to D(\mathcal{O}_{\Spec(k)}) = D(k)$
which we can identify with $K \mapsto R\Gamma(X, K)$. We also use
that the derived category $D(k)$ of $k$-modules is the same as the
category of graded $k$-vector spaces.

\medskip\noindent
Part (6) follows from Lemma \ref{lemma-dualizing-module-proper-over-A}
for coherent $\mathcal{F}$ and in general by unwinding
(5) for $K = \mathcal{F}[0]$ and $i = -\dim(X)$.

\medskip\noindent
Part (7) follows from Lemma \ref{lemma-smooth-proper}.
\end{proof}

\end{document}
